\documentclass[11pt]{article}
\usepackage{amsmath,amssymb}
\usepackage[margin=1in]{geometry}
\usepackage{booktabs}
\usepackage{hyperref}
\usepackage{xcolor}
\hypersetup{colorlinks=true, linkcolor=blue!50!black, urlcolor=blue!50!black}

\newcommand{\Psib}{\Psi}
\newcommand{\Var}{\mathcal{V}}
\newcommand{\code}[1]{\texttt{#1}}
\DeclareMathOperator{\Sol}{Sol}
\DeclareMathOperator{\ld}{ld}

\newtheorem{definition}{Definition}
\newtheorem{example}{Example}

\title{Membership versus variety, and the partial-solution strata:\\
       what changes if the PDE is added to the system\\ before
       decomposition}
\author{}
\date{\today}

\begin{document}
\maketitle

\noindent
A companion note to \emph{A New Method of Solving PDEs}. The core algorithm
\emph{reduces} the PDE $P$ modulo the ansatz $A$ and projects the remainder
onto the constants, returning the locus $\Var$ of constants for which the
ansatz family solves the PDE. A natural-looking alternative is to adjoin $P$
to the ansatz and run a differential decomposition (Rosenfeld--Gr\"obner or,
better, differential Thomas) on the \emph{combined} system $A \cup \{P\}$.
This note records why that answers a \emph{different} question, names the
discrepancy --- the \emph{partial-solution strata} --- and exhibits it on a
two-line example.

\section{Two different loci}

Let $A(c) = \{a_1, \dots, a_k\}$ be the ansatz, a differentially triangular
set whose coefficients involve constant parameters $c \in \mathbb{C}^n$, and
let $P$ be the PDE. The core method computes the \textbf{membership locus}
%
\begin{equation}
\label{eq:membership}
\Var \;=\; \{\, c \in \mathbb{C}^n : P \in [A(c)] \,\},
\end{equation}
%
a \emph{universal} condition: \emph{every} function in the ansatz family at
$c$ solves the PDE. Operationally, Ritt's full reduction of $P$ modulo $A$
yields a remainder $r = \sum_\alpha m_\alpha\, p_\alpha(c)$ with the $m_\alpha$
distinct power products in the non-constant indeterminates; $r$ vanishes
\emph{identically} iff every coefficient $p_\alpha(c)$ vanishes, so
$\Var = V(p_1, \dots, p_N)$.

Decomposing the combined system computes instead the \textbf{joint solution
variety}
%
\begin{equation}
\label{eq:variety}
\Sol(A) \cap \Sol(P),
\end{equation}
%
an \emph{existential} condition: there \emph{exists} a function satisfying both
$A$ and $P$. Projecting \eqref{eq:variety} onto the constants gives
%
\begin{equation}
\label{eq:overcount}
\pi_c\big(\Sol(A) \cap \Sol(P)\big)
   \;=\; \Var \;\cup\; \{\text{partial-solution strata}\}
   \;\supseteq\; \Var .
\end{equation}

\begin{definition}[partial-solution stratum]
\label{def:partial}
A \emph{partial-solution stratum} is a set of constant values $c$ at which a
\emph{proper, nontrivial} sub-family of the ansatz solutions satisfies $P$,
while the family as a whole does not. Equivalently: $c \notin \Var$, yet
$\Sol(A(c)) \cap \Sol(P)$ contains more than the zero solution.
\end{definition}

A prerequisite is immediate from Definition~\ref{def:partial}: the ansatz
solution space at $c$ must be at least two-dimensional. A one-dimensional
family either solves $P$ in its entirety or only through its zero member;
there is no proper nontrivial sub-family to peel off.

\section{A two-dimensional ansatz where the strata appear}

\begin{example}
\label{ex:partial}
Take a two-dimensional ansatz and a first-order PDE,
%
\begin{equation}
\label{eq:ex2}
A(c): \quad \Psib'' - c\,\Psib = 0,
\qquad\qquad
P: \quad \Psib' - \Psib = 0 .
\end{equation}
%
For $c = \omega^2$ the ansatz solution space is two-dimensional,
$\Psib = K_1 e^{\omega x} + K_2 e^{-\omega x}$.
\end{example}

\paragraph{Membership (core method): reduce $P$ modulo $A$.}
The leader of $P$ is $\Psib'$, which is ranked \emph{below} the leader $\Psib''$
of $A$, so $P$ is already fully reduced with respect to $A$; the remainder is
$P$ itself,
%
\[
r \;=\; \Psib' - \Psib \;=\; (+1)\cdot\Psib' + (-1)\cdot\Psib .
\]
%
The projection coefficients are the constants $+1$ and $-1$: never zero, with
no $c$-dependence. Hence
%
\[
\Var \;=\; V(1, -1) \;=\; \varnothing .
\]
%
This is correct: no value of $c$ makes the entire two-parameter family
$\{K_1 e^{\omega x} + K_2 e^{-\omega x}\}$ obey the first-order condition
$\Psib' = \Psib$.

\paragraph{Combined decomposition: adjoin $P$, rank the constant lowest.}
The decomposition selects the lower-leader equation $\Psib' - \Psib$ into the
triangular set first, then reduces $\Psib'' - c\,\Psib$ against it. Using
$\Psib' = \Psib \Rightarrow \Psib'' = \Psib$,
%
\[
\Psib'' - c\,\Psib \;\longrightarrow\; \Psib - c\,\Psib \;=\; (1-c)\,\Psib \;=\; 0 .
\]
%
This equation has leader $\Psib$ and initial $(1-c)$, so the algorithm splits
on the vanishing of $(1-c)$:
%
\begin{center}
\begin{tabular}{lll}
\toprule
cell & reduced system & solutions \\
\midrule
$c \ne 1$ & $(1-c)\Psib = 0,\ \ \Psib' = \Psib$ & $\Psib \equiv 0$ (trivial only)\\
$c = 1$   & $0 = 0,\ \ \Psib' = \Psib$           & $\Psib = K e^{x}$ (one-parameter)\\
\bottomrule
\end{tabular}
\end{center}

\noindent
The cell $c = 1$ is the partial-solution stratum. At $c = 1$ the ansatz has
the two modes $e^{x}$ and $e^{-x}$; the PDE $\Psib' = \Psib$ retains the
$e^{x}$ mode and discards $e^{-x}$. Exactly half the family solves the PDE ---
``some, but not all'' --- so $c = 1 \notin \Var$ even though it is a genuine,
nonempty solution stratum of the joint system.

\section{Matching the orders does not help}
\label{sec:orders}

Example~\ref{ex:partial} pairs a \emph{second}-order ansatz with a
\emph{first}-order PDE, which invites the hope that a partial stratum is an
order-mismatch artifact: raise the PDE to the ansatz's order and the gap might
close. It does not. The reason is structural. Ritt's full reduction of $P$
modulo $A$ rewrites $\Psib^{(n)}$ (and every higher derivative) through the
ansatz's order-$n$ ODE, so the \emph{remainder is always confined to order
$\le n-1$}, whatever the order of $P$. The PDE's content at and above the
ansatz leader is washed out; only the reduced remainder --- an operator of
order $\le n-1$ --- can distinguish the ansatz modes, and whether it shares a
\emph{proper} right factor with the ansatz operator is a condition on the
constants, indifferent to how $P$'s order compares with $n$. What closes the
gap is $n \le 1$ (a one-dimensional solution space; \S\ref{sec:onedim}),
a property of the \emph{ansatz}, not of the PDE or of the two orders' relation.

\begin{example}[both components second order]
\label{ex:both2}
%
\begin{equation}
\label{eq:both2}
A(c): \quad \Psib'' - c\,\Psib = 0,
\qquad\qquad
P: \quad \Psib'' - 3\Psib' + 2\Psib = 0 .
\end{equation}
%
Both are second order. The ansatz has the two-mode family
$\{e^{\sqrt{c}\,x}, e^{-\sqrt{c}\,x}\}$; the PDE has solutions $\{e^{x},
e^{2x}\}$.
\end{example}

\paragraph{Membership.}
Reducing $P$ modulo $A$ ($\Psib'' \equiv c\,\Psib$, with multiplier $h = 1$,
so $h \neq 0$ \emph{everywhere} --- the strata below are genuine, not
bad-locus artifacts) gives
\[
r \;=\; c\,\Psib - 3\Psib' + 2\Psib \;=\; -3\,\Psib' + (c+2)\,\Psib .
\]
The projection coefficients are $-3$ and $c+2$; the first never vanishes, so
$\Var = \varnothing$.

\paragraph{Existence.}
A mode $e^{\rho x}$ ($\rho = \pm\sqrt{c}$) solves $P$ iff
$\rho^2 - 3\rho + 2 = 0$, i.e.\ $\rho \in \{1, 2\}$, i.e.\ $c = \rho^2 \in
\{1, 4\}$:
%
\begin{center}
\begin{tabular}{llll}
\toprule
$c$ & ansatz modes & solves $P$ & other mode \\
\midrule
$1$ & $e^{\pm x}$  & $e^{x}$  ($1-3+2=0$) & $e^{-x}$: $1+3+2 = 6 \neq 0$ \\
$4$ & $e^{\pm 2x}$ & $e^{2x}$ ($4-6+2=0$) & $e^{-2x}$: $4+6+2 = 12 \neq 0$ \\
\bottomrule
\end{tabular}
\end{center}
%
\noindent
So $V_{\exists\backslash\Psib} = \{1, 4\}$ while $\Var = \varnothing$: at each
of $c = 1, 4$ exactly one of the two ansatz modes survives $P$, a genuine
partial stratum with $h \neq 0$. The GCRD test agrees: the reduced remainder
is $R \sim \partial - \tfrac{c+2}{3}$, which right-divides $\partial^2 - c$ iff
$\big(\tfrac{c+2}{3}\big)^2 = c$, i.e.\ $c^2 - 5c + 4 = (c-1)(c-4) = 0$, the
order-$1$ (partial-stratum) case of the trichotomy. The example is
Example~\ref{ex:partial} with its first-order PDE replaced by a second-order
one that reduces to the same kind of order-$1$ remainder; the two behave
identically precisely because reduction erases $P$'s second-order part.

\section{The correct selection criterion}

Equation~\eqref{eq:overcount} warns that projecting the \emph{whole} combined
variety over-counts $\Var$. One might hope to recover $\Var$ by keeping the
cell in which the adjoined PDE-equation collapsed to $0 = 0$; Example
\ref{ex:partial} shows this is wrong, as that cell is $c = 1$ whereas
$\Var = \varnothing$. The collapse cell is the \emph{existential} locus, not
the membership locus.

The correct criterion is \textbf{dimension preservation}: $c \in \Var$ iff the
combined solution space at $c$ still equals the \emph{full} ansatz solution
space, i.e. the PDE removed nothing. In Example~\ref{ex:partial} no cell
retains the full two-dimensional space (the $c=1$ cell keeps only a
one-dimensional slice), so $\Var = \varnothing$, consistent with the direct
reduction.

\section{Contrast: a one-dimensional ansatz has no partial strata}
\label{sec:onedim}

Shrink the ansatz to one dimension and test a second-order PDE:
%
\[
A(c): \quad \Psib' - c\,\Psib = 0 \quad(\Psib = K e^{cx}),
\qquad\qquad
P: \quad \Psib'' - \Psib = 0 .
\]
%
Now $P$'s leader $\Psib''$ dominates $A$'s leader $\Psib'$. Reducing $P$ modulo
$A$ gives $\Psib'' \to c^2 \Psib$, hence $r = (c^2 - 1)\,\Psib$ and
%
\[
\Var \;=\; V(c^2 - 1) \;=\; \{\,+1,\,-1\,\} .
\]
%
The combined decomposition produces the cells
%
\[
\{\, c^2 \ne 1 : \Psib \equiv 0 \,\}, \qquad
\{\, c = 1   : \Psib = K e^{x} \,\}, \qquad
\{\, c = -1  : \Psib = K e^{-x} \,\},
\]
%
and the nontrivial cells $c = \pm 1$ now retain the \emph{entire}
(one-dimensional) family. Here the collapse cells coincide with $\Var$ and
there are no partial strata: a one-dimensional family admits no proper
nontrivial sub-family.

\section{Autoreducedness does not preclude the strata: a PDE example}

One might conjecture a combinatorial precondition for the partial strata:
that they arise only when the reduced PDE and the ansatz still \emph{interact}
under Ritt reduction. In Example~\ref{ex:partial} the remainder
$r = \Psib' - \Psib$ can reduce the ansatz equation $\Psib'' - c\,\Psib$ (its
leader $\Psib'$ has $\Psib''$ as a proper derivative), and it was exactly that
rewriting which produced the split on $1 - c$. What if reduction terminates
completely --- $r$ can reduce nothing in $A$ and $A$ nothing in $r$, so that
$A \cup \{r\}$ is autoreduced? Does that preclude a partial-solution stratum?

In the ordinary-differential setting (one derivation, one unknown) the
conjecture holds, but only vacuously. There $A$ is a single equation with
leader $\Psib^{(k)}$, and the reduced remainder involves only
$\Psib, \dots, \Psib^{(k-1)}$; if $r$ involves the unknown at all, its leader
$\Psib^{(j)}$ has $\Psib^{(k)}$ as a proper derivative, so $r$ \emph{can}
always reduce $A$. The hypothesis therefore forces $r = p(c)$, a pure
constants condition --- and then both computations return $V(p)$ carrying the
full ansatz family: dimension-preserving, no strata.

With several derivations, leaders can be incomparable and the conjecture
fails.

\begin{example}
\label{ex:pde}
Two derivations $\partial_x, \partial_y$:
%
\[
A(c): \quad \Psib_{yy} - c\,\Psib = 0,
\qquad\qquad
P: \quad \Psib_x - \Psib = 0 .
\]
\end{example}

\noindent
$P$ is already reduced with respect to $A$ ($\Psib_x$ is not a derivative of
$\Psib_{yy}$), so $r = P$; and $r$ can reduce nothing in $A$, whose equation
contains only $\Psib_{yy}$ and $\Psib$, neither a derivative of $\Psib_x$.
Thus $A \cup \{r\}$ is autoreduced --- and even \emph{coherent}: the leaders'
least common derivative is $\Psib_{xyy}$, and the $\Delta$-polynomial
%
\[
\partial_y^2\, r - \partial_x\, A
\;=\; (\Psib_{xyy} - \Psib_{yy}) - (\Psib_{xyy} - c\,\Psib_x)
\;=\; c\,\Psib_x - \Psib_{yy}
\;\longrightarrow\; c\,\Psib - c\,\Psib \;=\; 0 .
\]
%
The combined decomposition therefore emits $A \cup \{r\}$ unchanged, as a
single cell, with no splitting at all.

Membership: the projection coefficients of $r$ are $+1$ and $-1$, so
$\Var = \varnothing$ --- correct, since the general ansatz solution
$\Psib = F(x)\,e^{\omega y} + G(x)\,e^{-\omega y}$ ($\omega^2 = c$; $F$, $G$
arbitrary functions) does not satisfy $\Psib_x = \Psib$.

Joint variety: $\Psib_x = \Psib$ forces $\Psib = e^{x} h(y)$, and then
$h'' = c\,h$, so
%
\[
\Psib \;=\; e^{x}\big(K_1 e^{\omega y} + K_2 e^{-\omega y}\big)
\]
%
--- nontrivial at \emph{every} $c$, and a proper subfamily of the ansatz
family (the arbitrary functions $F$, $G$ collapse to constant multiples of
$e^{x}$). Every $c$ is a partial-solution stratum; here the stratum is the
whole parameter space.

The structural reading: ``$A \cup \{r\}$ is autoreduced and coherent'' means
the decomposition has no more algebra to do --- it is precisely a finished
cell --- \emph{not} that $r$ imposes no constraint. The cell's solution set
is the joint variety, and whenever that sits strictly between $0$ and
$\Sol(A(c))$ at some $c \notin \Var$, a partial stratum is present. What
precludes the strata is $r$ having no non-constant part, $r = p(c)$ ---
which, in the single-derivation case, is the only way the autoreducedness
hypothesis can be met at all. The strata are a fact of the solution
geometry, not of the rewriting combinatorics.

\section{Reading}

The membership formulation~\eqref{eq:membership} asks: \emph{for which
ODE-coefficients $c$ is the PDE an automatic consequence of the ansatz?} ---
so that the ansatz family is, as a whole, a PDE-solution family. The partial
strata answer a different question: \emph{for which $c$ does the ansatz family
happen to contain a PDE solution it was not designed to?} --- coincidental
solutions living on a proper sub-family at special parameter values. They are
real solutions, but invisible to the membership formulation by construction.

For the hydrogen computation this means: adjoining the Schr\"odinger equation
to the ansatz and decomposing surfaces parameter values where a degenerate,
lower-dimensional slice of the ansatz solves the PDE while the generic member
does not. If the goal is the membership locus $\Var$, either reduce $P$ modulo
$A$ directly, or --- having decomposed the combined system --- keep only the
dimension-preserving cells. Projecting the full joint variety answers the
existential question and returns $\Var$ together with the partial strata.

\end{document}
